\documentclass{article}
\usepackage{amsmath,amssymb,amsthm}
\usepackage{algorithm,algorithmic}
\usepackage{hyperref}
\usepackage{xcolor}
\newtheorem{theorem}{Theorem}
\newtheorem{lemma}{Lemma}
\newcommand{\EE}{\mathbb{E}}
\newcommand{\RR}{\mathbb{R}}
\newcommand{\norm}[1]{\left\|#1\right\|}
\begin{document}

\section*{Adaptive Variance‑Reduced Gradient (AVRG)}

\section*{Mathematical Formulation}
Let $f:\RR^{d}\rightarrow\RR$ be $L$–smooth and satisfy the $\mu$–Polyak–Łojasiewicz (PL) condition.  
The algorithm maintains a moving‑average estimate of the gradient variance

\[
\boxed{v_{t}= (1-\alpha)v_{t-1}+ \alpha \nabla f(x_{t})},\qquad \alpha\in(0,1] .
\tag{1}
\]

Because $v_{t}$ is biased toward zero at early iterations, a bias‑corrected version is used

\[
\boxed{\hat v_{t}= \frac{v_{t}}{1-(1-\alpha)^{t}}}.
\tag{2}
\]

The step size is adapted to the magnitude of $\hat v_{t}$:

\[
\boxed{\eta_{t}= \frac{\eta}{\,1+\beta \|\hat v_{t}\|\,}},\qquad \eta>0,\;\beta\ge 0 .
\tag{3}
\]

Given a stochastic gradient estimator $g_{t}$ (e.g. SARAH/SPIDER), the iterate update is

\[
\boxed{x_{t+1}= x_{t}-\eta_{t} g_{t}} .
\tag{4}
\]

When the full gradient is used ($g_{t}=\nabla f(x_{t})$) the algorithm reduces to a deterministic adaptive method.

\section*{Pseudocode}
\begin{algorithm}[H]
\caption{Adaptive Variance‑Reduced Gradient (AVRG)}
\label{alg:avrg}
\begin{algorithmic}[1]
\STATE \textbf{Input:} initial point $x_{0}$, step size $\eta$, decay $\beta$, variance decay $\alpha$, batch size $b$ (for SARAH/SPIDER), max iterations $T$.
\STATE Initialize $v_{0}=0$, $g_{0}= \frac{1}{b}\sum_{i\in\mathcal{B}_{0}}\nabla f_{i}(x_{0})$.
\FOR{$t=0,\dots,T-1$}
    \STATE Compute a stochastic gradient $g_{t}$ (e.g. SARAH update)
    \STATE Update variance estimator: $v_{t+1}= (1-\alpha)v_{t}+ \alpha g_{t}$ \hfill\textcolor{gray}{\# (1)}
    \STATE Bias correction: $\hat v_{t+1}= v_{t+1}/\bigl(1-(1-\alpha)^{t+1}\bigr)$ \hfill\textcolor{gray}{\# (2)}
    \STATE Adaptive step size: $\displaystyle \eta_{t}= \frac{\eta}{1+\beta \|\hat v_{t+1}\|}$ \hfill\textcolor{gray}{\# (3)}
    \STATE Update iterate: $x_{t+1}= x_{t}-\eta_{t} g_{t}$ \hfill\textcolor{gray}{\# (4)}
\ENDFOR
\STATE \textbf{Return} $x_{T}$
\end{algorithmic}
\end{algorithm}

\section*{Dry Run Example}
Consider the scalar quadratic $f(w)=(w-3)^{2}$, thus $\nabla f(w)=2(w-3)$, $L=2$, $\mu=2$ (PL holds with equality).  
Set $\eta=0.1$, $\alpha=0.5$, $\beta=0$, $x_{0}=0$, and use the \emph{full} gradient ($g_{t}=\nabla f(x_{t})$).

\begin{center}
\begin{tabular}{c|c|c|c|c}
$t$ & $x_{t}$ & $g_{t}=2(x_{t}-3)$ & $v_{t}$ (by (1)) & $f(x_{t})$\\\hline
0 & $0$ & $-6$ & $v_{0}=0$ & $9$\\
1 & $x_{1}=0-0.1(-6)=0.6$ & $2(0.6-3)=-4.8$ & $v_{1}=0.5\cdot0+0.5(-6)=-3$ & $5.76$\\
2 & $x_{2}=0.6-0.1(-4.8)=1.08$ & $2(1.08-3)=-3.84$ & $v_{2}=0.5(-3)+0.5(-4.8)=-3.9$ & $3.6864$\\
3 & $x_{3}=1.08-0.1(-3.84)=1.464$ & $2(1.464-3)=-3.072$ & $v_{3}=0.5(-3.9)+0.5(-3.84)=-3.87$ & $2.354$\\
\end{tabular}
\end{center}
The objective values decrease monotonically, illustrating the algorithm’s descent property.

\newpage
\section*{Assumptions}
\begin{enumerate}
\item \textbf{Smoothness.} $f$ is $L$‑smooth:
\[
\forall x,y\in\RR^{d},\qquad 
f(y)\le f(x)+\langle\nabla f(x),y-x\rangle+\frac{L}{2}\|y-x\|^{2}.
\tag{A1}
\]
\item \textbf{Polyak–Łojasiewicz (PL) condition.}
\[
\forall x\in\RR^{d},\qquad 
\frac{1}{2}\|\nabla f(x)\|^{2}\ge \mu\bigl(f(x)-f^{\star}\bigr),
\tag{A2}
\]
where $f^{\star}= \min_{x}f(x)$ and $\mu>0$.
\item \textbf{Bounded variance of the stochastic estimator.} For the SARAH/SPIDER estimator $g_{t}$,
\[
\EE\!\left[\;g_{t}\mid \mathcal{F}_{t}\right]=\nabla f(x_{t}),\qquad 
\EE\!\left[\;\|g_{t}-\nabla f(x_{t})\|^{2}\mid\mathcal{F}_{t}\right]\le\sigma^{2},
\tag{A3}
\]
where $\mathcal{F}_{t}$ is the $\sigma$‑algebra generated by all randomness up to iteration $t$.
\item \textbf{Hyper‑parameter constraints.}
\[
0<\alpha\le1,\qquad \eta>0,\qquad \beta\ge0,\qquad 
\eta\le\frac{1}{L}.
\tag{A4}
\]
\end{enumerate}

\section*{Main Convergence Theorem}
\begin{theorem}[Linear convergence under PL]\label{thm:linear}
Assume (A1)–(A4) and let $\{x_{t}\}$ be generated by AVRG with the full‑gradient choice $g_{t}=\nabla f(x_{t})$.  
Define the maximal possible adaptive stepsize
\[
\bar\eta\;:=\;\frac{\eta}{1+\beta M},\qquad 
M\;:=\;\sup_{t\ge0}\norm{\hat v_{t}}.
\]
If $0<\bar\eta\le \frac{2}{L}$, then for all $t\ge0$
\[
\EE\!\bigl[f(x_{t})-f^{\star}\bigr]\;\le\;
\bigl(1-\rho\bigr)^{t}\bigl(f(x_{0})-f^{\star}\bigr),
\qquad 
\rho:=2\mu\bar\eta\Bigl(1-\frac{L\bar\eta}{2}\Bigr)\in(0,1).
\tag{5}
\]
Consequently, $\{x_{t}\}$ converges linearly to the unique minimiser $x^{\star}$.
\end{theorem}

\section*{Theoretical Properties}
\begin{itemize}
\item \textbf{Stability.} The adaptive stepsize $\eta_{t}$ is always bounded between $\eta/(1+\beta M)$ and $\eta$, guaranteeing that the descent Lemma (A1) can be applied at every iteration.
\item \textbf{Hyper‑parameter sensitivity.} Larger $\beta$ yields a smaller $\bar\eta$ and thus a slower but more stable convergence; the rate $\rho$ is monotone increasing in $\bar\eta$ up to the limit $\bar\eta=2/L$.
\item \textbf{Comparison to baselines.}
    \begin{itemize}
        \item With $\beta=0$ the algorithm reduces to a constant‑stepsize gradient method and recovers the classic PL linear rate $\rho=2\mu\eta(1-L\eta/2)$.
        \item When $\alpha\to1$, $v_{t}\approx g_{t}$, so $\|\hat v_{t}\|$ closely tracks $\|\nabla f(x_{t})\|$ and the stepsize automatically shrinks near the optimum, mimicking an ``Armijo’’ line‑search without extra function evaluations.
    \end{itemize}
\end{itemize}

\section*{Discussion}
The theorem shows that the moving‑average variance estimator does not deteriorate the PL‑induced linear convergence; it merely rescales the effective stepsize. The bias‑correction (2) guarantees that $M$ is finite even when $\alpha$ is small, because $\hat v_{t}\to \nabla f(x_{t})$ as $t\to\infty$. The same analysis extends to stochastic SARAH/SPIDER estimators by adding a standard variance term, which only affects the constant factor in the linear rate (see Lemma~\ref{lem:stoch}).

\newpage
\section*{Preliminaries (Definitions and Lemmas)}
\begin{lemma}[Descent Lemma]\label{lem:descent}
If $f$ is $L$‑smooth, then for any $x$ and any vector $d$
\[
f(x-d)\le f(x)-\langle\nabla f(x),d\rangle+\frac{L}{2}\|d\|^{2}.
\tag{L1}
\]
\end{lemma}
\begin{proof}
Directly from (A1) with $y=x-d$.
\end{proof}

\begin{lemma}[Bound on adaptive stepsize]\label{lem:stepsize}
Let $M:=\sup_{t}\|\hat v_{t}\|<\infty$. Then for all $t$
\[
\frac{\eta}{1+\beta M}\;\le\;\eta_{t}\;\le\;\eta .
\tag{L2}
\]
\end{lemma}
\begin{proof}
From definition (3), $\|\hat v_{t}\|\le M$ implies $1+\beta\|\hat v_{t}\|\le1+\beta M$, yielding the lower bound. The numerator $\eta$ gives the obvious upper bound.
\end{proof}

\begin{lemma}[Variance bound for SARAH/SPIDER]\label{lem:stoch}
Under (A3) and for any stepsize $\eta_{t}\le 2/L$, the iterates satisfy
\[
\EE\!\bigl[f(x_{t+1})-f^{\star}\mid\mathcal{F}_{t}\bigr]
\le f(x_{t})-f^{\star}
-\eta_{t}\Bigl(1-\frac{L\eta_{t}}{2}\Bigr)\|\nabla f(x_{t})\|^{2}
+\frac{L\eta_{t}^{2}\sigma^{2}}{2}.
\tag{L3}
\]
\end{lemma}
\begin{proof}
Apply Lemma~\ref{lem:descent} with $d=\eta_{t}g_{t}$, take conditional expectation, use unbiasedness and the variance bound (A3), and rearrange the terms.
\end{proof}

\section*{Main Proof (Step‑by‑Step Derivation)}
We prove Theorem~\ref{thm:linear} for the deterministic case $g_{t}=\nabla f(x_{t})$.  
All non‑trivial steps are labelled \textbf{[Step N]}.

\subsection*{Step 1 – Apply Descent Lemma}
Using Lemma~\ref{lem:descent} with $d=\eta_{t}\nabla f(x_{t})$,
\[
f(x_{t+1})\le f(x_{t})-\eta_{t}\|\nabla f(x_{t})\|^{2}
+\frac{L\eta_{t}^{2}}{2}\|\nabla f(x_{t})\|^{2}.
\tag{6}
\]
\textbf{[Step 1]}

\subsection*{Step 2 – Rearrange}
Factor $\|\nabla f(x_{t})\|^{2}$:
\[
f(x_{t+1})\le f(x_{t})-
\eta_{t}\Bigl(1-\frac{L\eta_{t}}{2}\Bigr)\|\nabla f(x_{t})\|^{2}.
\tag{7}
\]
\textbf{[Step 2]}

\subsection*{Step 3 – Insert PL inequality}
From (A2),
\[
\|\nabla f(x_{t})\|^{2}\ge 2\mu\bigl(f(x_{t})-f^{\star}\bigr).
\tag{8}
\]
\textbf{[Step 3]}

\subsection*{Step 4 – Combine (7) and (8)}
\[
f(x_{t+1})-f^{\star}
\le\Bigl[1-2\mu\eta_{t}\Bigl(1-\frac{L\eta_{t}}{2}\Bigr)\Bigr]\bigl(f(x_{t})-f^{\star}\bigr).
\tag{9}
\]
\textbf{[Step 4]}

\subsection*{Step 5 – Replace $\eta_{t}$ by its lower bound}
From Lemma~\ref{lem:stepsize}, $\eta_{t}\ge\bar\eta:=\eta/(1+\beta M)$. Define
\[
\rho\;:=\;2\mu\bar\eta\Bigl(1-\frac{L\bar\eta}{2}\Bigr).
\tag{10}
\]
Because $0<\bar\eta\le 2/L$, we have $0<\rho<1$. Using $\eta_{t}\ge\bar\eta$ in (9) yields
\[
f(x_{t+1})-f^{\star}
\le (1-\rho)\bigl(f(x_{t})-f^{\star}\bigr).
\tag{11}
\]
\textbf{[Step 5]}

\subsection*{Step 6 – Unroll the recursion}
Iterating (11) from $0$ to $t-1$ gives
\[
f(x_{t})-f^{\star}
\le (1-\rho)^{t}\bigl(f(x_{0})-f^{\star}\bigr).
\tag{12}
\]
\textbf{[Step 6]}

\subsection*{Step 7 – Take expectations}
Since the deterministic case involves no randomness, $\EE[\cdot]=\cdot$, and (12) is exactly the claim (5).

Thus Theorem~\ref{thm:linear} holds. \hfill $\blacksquare$

\section*{Rate Derivation (With Constants)}
The contraction factor $\rho$ can be expressed explicitly in terms of the original hyper‑parameters:
\[
\rho
= \frac{2\mu\eta}{1+\beta M}
\Bigl(1-\frac{L\eta}{2(1+\beta M)}\Bigr).
\tag{13}
\]
If we set $\beta=0$, then $M$ disappears and we recover the classical PL rate
$\rho_{\text{GD}}=2\mu\eta(1-L\eta/2)$.  
Increasing $\beta$ reduces $\rho$ proportionally, reflecting the more conservative stepsizes.

\section*{Special Cases}
\begin{enumerate}
\item \textbf{Constant stepsize (β=0).} The algorithm coincides with plain gradient descent; Theorem~\ref{thm:linear} reduces to the standard PL linear convergence.
\item \textbf{Full variance tracking (α=1).} Then $v_{t}=g_{t}$ and $\hat v_{t}=g_{t}$, so $M=\sup_{t}\|\nabla f(x_{t})\|$, a finite quantity for $L$‑smooth functions on a bounded level set. The adaptive stepsize automatically scales with the true gradient norm.
\item \textbf{Stochastic SARAH/SPIDER estimator.} Using Lemma~\ref{lem:stoch} instead of (6) adds the variance term $\frac{L\eta_{t}^{2}\sigma^{2}}{2}$, which yields
\[
\EE[f(x_{t+1})-f^{\star}]
\le (1-\rho)\EE[f(x_{t})-f^{\star}]
+\frac{L\eta^{2}\sigma^{2}}{2(1+\beta M)^{2}} .
\]
Hence the iterates converge linearly to a neighbourhood of size $O(\eta^{2}\sigma^{2})$.
\end{enumerate}

\end{document}